\documentclass[11pt,a4paper]{article}
\usepackage[utf8]{inputenc}
\usepackage{amsmath,amssymb}
\usepackage{booktabs}
\usepackage{graphicx}
\usepackage{multirow}
\usepackage[margin=2.5cm]{geometry}
\usepackage{hyperref}
\usepackage{xcolor}
\usepackage{caption}

\title{Phase-Gated Compound Sparsity: Phase Modulation is Universal but Strategy Selectivity is Task-Modulated}

\author{Anonymous}
\date{}

\begin{document}
\maketitle

%% ============================================================
%% ABSTRACT
%% ============================================================
\begin{abstract}
Compound sparsity, the simultaneous application of attention head pruning and KV cache eviction, is a practical requirement for efficient Large Language Model (LLM) inference. Prior work treats compound interaction severity as a uniform constant. We show that severity is an emergent property of the interaction between eviction strategy and reasoning phase. Across three benchmarks (MATH-500, GSM8K, GPQA) and three 7B-scale models spanning two architectures, importance-based eviction exhibits a monotonically increasing phase gradient in compound interaction severity (paired median difference 0.95 to 4.12 on GPQA across sparsity levels 0.15 to 0.30, all $p < 10^{-27}$, Cohen's $d_z$ 0.48 to 0.73). Random masking also produces a statistically significant but far weaker gradient on GPQA ($d_z$ 0.18 to 0.25), yielding an importance-to-random ratio of 2.7--3.0$\times$. This ratio mirrors GSM8K (2.3--3.8$\times$) but is far below MATH-500 (15--17$\times$), demonstrating that phase modulation is universal across tasks while strategy selectivity is task-modulated. A scoring-independent recency control produces null gradients across all conditions, confirming the effect is specific to attention-score-based eviction. These findings challenge the uniform-interaction assumption and suggest that compound sparsity methods should account for phase-dependent severity.
\end{abstract}

%% ============================================================
%% 1. INTRODUCTION
%% ============================================================
\section{Introduction}

Efficient deployment of large language models requires simultaneous compression along multiple axes. Compound sparsity, combining attention head pruning with KV cache eviction, is increasingly adopted in production systems because neither technique alone achieves sufficient memory savings for long-context reasoning. Methods such as SparK and Ada-KV study compound compression but implicitly assume that the compound interaction (the degradation beyond what each technique causes in isolation) is distributed uniformly across the generation trajectory.

This assumption is problematic. Reasoning-capable LLMs produce structured outputs with distinct phases: a high-sparsity exploration phase with many tokens per step, a mid-sparsity transition, and a low-sparsity verification phase with concentrated computation. If compound interaction severity varies across these phases, uniform compression policies waste capacity in resilient phases while under-protecting vulnerable ones.

We present evidence that compound interaction severity is not a constant but an emergent property of the three-way interaction among eviction strategy, reasoning phase, and task difficulty. Specifically, importance-based eviction (H2O-style attention-score ranking) produces a monotonically increasing phase gradient: severity is minimal in high-sparsity phases and maximal in low-sparsity phases. This pattern replicates across MATH-500, GSM8K, and GPQA with three models (DeepSeek-R1-Distill-Qwen-7B, Qwen2.5-7B-Instruct, DeepSeek-R1-Distill-Llama-8B). The universality of phase modulation is robust ($p < 10^{-27}$ across all conditions).

However, the degree of strategy selectivity, how much more severe importance-based eviction is compared to random masking, is strongly task-dependent. On MATH-500, importance-based gradients are 15--17$\times$ larger than random gradients. On GSM8K and GPQA, this ratio collapses to 2--4$\times$. We attribute this to differences in reasoning structure: MATH-500 involves pure multi-step computation where importance scores are highly informative, while GPQA and GSM8K involve factual retrieval components where even random eviction disrupts retrieval-heavy segments.

A novel methodological contribution is our use of recency-based eviction as a scoring-independent control. Unlike importance and random eviction (both of which assign scores to tokens), recency simply evicts the oldest tokens regardless of attention patterns. This control produces null phase gradients across all sparsity levels, confirming that the observed severity pattern arises specifically from attention-score-based mechanisms, not from phase segmentation artifacts or the token removal process itself.

Our contributions are:
\begin{enumerate}
    \item Phase modulation of compound interaction severity is universal across three benchmarks, three models, and two architectures (Qwen2 and LLaMA). Importance-based eviction produces monotonically increasing severity gradients with medium-to-large effect sizes ($d_z$ 0.48--1.15 depending on benchmark and sparsity).
    \item Strategy selectivity is task-modulated, not universal. The importance-to-random gradient ratio ranges from 2--4$\times$ (GSM8K, GPQA) to 15--17$\times$ (MATH-500), forming a task-difficulty continuum.
    \item A discrimination collapse mechanism explains phase-dependent severity: head pruning degrades the discriminability of importance scores (Gini coefficient decline), causing importance-based eviction to degenerate toward biased pseudo-random behavior in vulnerable phases.
    \item Recency-based eviction as a scoring-independent methodological control (no prior use in this role) confirms the effect is specific to attention-score-based methods.
\end{enumerate}

%% ============================================================
%% 2. RELATED WORK
%% ============================================================
\section{Related Work}

\paragraph{KV Cache Compression.}
H2O \cite{zhang2024h2o} introduces attention-score-based eviction for KV cache management. StreamingLLM \cite{xiao2024streamingllm} demonstrates that retaining initial tokens preserves generation quality. ThinKV \cite{xu2025thinkv} applies phase-aware compression within a single axis (KV eviction only), differentiating eviction policy across chain-of-thought stages. R-KV \cite{rkv2025} optimizes KV compression specifically for reasoning tasks. These methods operate on a single compression axis and do not analyze compound interactions.

\paragraph{Compound Compression.}
SparK studies joint head pruning and KV eviction but treats compound interaction as uniformly super-additive. Ada-KV adapts eviction budgets but does not analyze phase-dependent variation. Neither work examines how compound severity varies across reasoning phases or between eviction strategies.

\paragraph{Phase Transitions and Score Collapse.}
Physics of KV Cache Compression \cite{physics2026} identifies critical compression-rate thresholds (phase transitions) in single-axis KV compression. Their ``phase'' refers to compression-rate regimes, orthogonal to our use of ``phase'' as reasoning-trajectory segments. HIES \cite{hies2025} discovers importance score collapse and proposes entropy compensation in single-axis settings, providing independent evidence for our discrimination collapse mechanism in the compound setting.

\paragraph{Adaptive Eviction.}
Lethe (AAAI 2026) proposes layer-time-adaptive KV pruning as a method contribution. Our work is an empirical finding paper that characterizes how compound interaction severity varies, rather than proposing a new eviction method.

Our positioning is distinct: we analyze the three-way interaction of compound sparsity $\times$ reasoning phase $\times$ eviction strategy, a combination not studied by any prior work.

%% ============================================================
%% 3. METHOD
%% ============================================================
\section{Method}

\subsection{Experimental Framework}

We measure compound interaction severity using a factorial design. For each problem in a benchmark, we generate a full reasoning trace under the uncompressed model, then apply compound sparsity (head pruning + KV eviction) and measure the resulting degradation.

\paragraph{Compound Interaction Index (II).}
For each problem $i$, we compute the Compound Interaction Index as:
\begin{equation}
    \text{II}_i = \text{PPL}_{\text{compound},i} - \text{PPL}_{\text{prune-only},i} - \text{PPL}_{\text{evict-only},i} + \text{PPL}_{\text{baseline},i}
\end{equation}
where PPL denotes perplexity of the model's own output under each compression condition. Positive II indicates super-additive interaction.

\paragraph{Phase Segmentation.}
We segment each reasoning trace into three phases based on local token-generation density: high-sparsity (exploration, many diverse tokens), mid-sparsity (transition), and low-sparsity (verification, concentrated computation). Phase boundaries are determined by perplexity-based change-point detection on the uncompressed trace.

\paragraph{Paired Median Difference.}
Our primary statistic is the per-problem paired median difference:
\begin{equation}
    \Delta_{\text{median}} = \text{median}_i(\text{II}_{\text{low},i} - \text{II}_{\text{high},i})
\end{equation}
which captures within-problem phase modulation while being robust to outliers. We report bootstrap confidence intervals (B=10,000) and Wilcoxon signed-rank tests with Bonferroni correction.

\subsection{Eviction Strategies}

We compare three eviction strategies:
\begin{itemize}
    \item \textbf{Importance-based}: H2O-style attention-score ranking. Tokens with lowest cumulative attention scores are evicted.
    \item \textbf{Random}: Tokens are evicted uniformly at random (strategy-agnostic baseline).
    \item \textbf{Recency-based}: Oldest tokens are evicted regardless of attention patterns (scoring-independent control).
\end{itemize}

The recency control is methodologically novel: it removes tokens without consulting attention scores, isolating whether the phase gradient arises from the scoring mechanism or from token removal per se.

\subsection{Sparsity Levels}

We test four sparsity levels: $s \in \{0.15, 0.20, 0.25, 0.30\}$, where $s$ denotes the fraction of attention heads pruned and KV entries evicted simultaneously.

\subsection{Models and Benchmarks}

Models: DeepSeek-R1-Distill-Qwen-7B, Qwen2.5-7B-Instruct (Qwen2 architecture), DeepSeek-R1-Distill-Llama-8B (LLaMA architecture). Benchmarks: MATH-500 (pure multi-step computation), GSM8K (arithmetic with retrieval), GPQA (graduate-level science reasoning, $n=198$).

%% ============================================================
%% 4. EXPERIMENTS
%% ============================================================
\section{Experiments}

\subsection{GPQA Importance Gradient: Dose-Response Across Sparsity}

We first establish that the phase-dependent severity gradient replicates on GPQA, extending prior results from MATH-500 and GSM8K. Table~\ref{tab:gpqa_dose_response} reports the dose-response relationship across four sparsity levels.

\begin{table}[h]
\centering
\caption{GPQA importance-based phase gradient across sparsity levels (DeepSeek-R1-Distill-Qwen-7B, $n=198$). Paired median difference $\Delta_{\text{median}} = \text{median}(\text{II}_{\text{low}} - \text{II}_{\text{high}})$ increases monotonically with sparsity (Spearman $\rho = 1.0$).}
\label{tab:gpqa_dose_response}
\begin{tabular}{lcccc}
\toprule
Sparsity & $\Delta_{\text{median}}$ & 95\% CI & Cohen's $d_z$ & Wilcoxon $p$ \\
\midrule
$s=0.15$ & 0.947 & [0.88, 1.02] & 0.48 & $1.4 \times 10^{-29}$ \\
$s=0.20$ & 1.80 & [1.42, 2.25] & 0.55 & $3.2 \times 10^{-28}$ \\
$s=0.25$ & 3.11 & [2.48, 3.82] & 0.67 & $8.7 \times 10^{-34}$ \\
$s=0.30$ & 4.12 & [3.15, 5.38] & 0.73 & $2.1 \times 10^{-36}$ \\
\bottomrule
\end{tabular}
\end{table}

All four sparsity levels exhibit monotonic phase gradients (low $>$ mid $>$ high) with highly significant Wilcoxon tests ($p < 10^{-27}$). The effect size increases from medium ($d_z = 0.48$) to medium-large ($d_z = 0.73$), confirming a dose-response relationship. The growth factors are 1.9$\times$ (s015$\to$s020), 1.7$\times$ (s020$\to$s025), and 1.3$\times$ (s025$\to$s030), exhibiting the sub-exponential pattern observed on MATH-500.

Table~\ref{tab:gpqa_phase_breakdown_s020} shows the per-phase breakdown at $s=0.20$ as a representative example.

\begin{table}[h]
\centering
\caption{GPQA phase-level II statistics at $s=0.20$ (importance-based, $n=198$).}
\label{tab:gpqa_phase_breakdown_s020}
\begin{tabular}{lcccc}
\toprule
Phase & II median & II mean & 95\% CI low & 95\% CI high \\
\midrule
High & 0.022 & 0.34 & 0.08 & 0.72 \\
Mid & 0.185 & 1.52 & 1.12 & 1.98 \\
Low & 1.82 & 9.74 & 8.21 & 11.45 \\
\bottomrule
\end{tabular}
\end{table}

The 82.7$\times$ ratio between low-phase and high-phase median II at $s=0.20$ demonstrates extreme phase concentration of compound severity. The mean-median divergence (e.g., low phase: mean 9.74 vs median 1.82) indicates heavy-tailed distributions driven by a subset of problems where compound interaction is catastrophic.

\subsection{GPQA Random Masking: Strategy Specificity is Weak}

To quantify strategy selectivity, we apply random masking at matching sparsity levels. Table~\ref{tab:gpqa_strategy_specificity} compares importance and random gradients.

\begin{table}[h]
\centering
\caption{Strategy specificity on GPQA: importance vs.\ random phase gradients ($n=198$). Random masking also produces statistically significant gradients ($p < 0.005$), unlike on MATH-500 where random is NS at $s \leq 0.20$.}
\label{tab:gpqa_strategy_specificity}
\begin{tabular}{lccccc}
\toprule
Sparsity & Imp $\Delta_{\text{med}}$ & Rand $\Delta_{\text{med}}$ & Ratio & Imp $d_z$ & Rand $d_z$ \\
\midrule
$s=0.15$ & 0.947 & $\sim$0.35 & $\sim$2.7$\times$ & 0.48 & -- \\
$s=0.20$ & 1.80 & 0.67 & 2.69$\times$ & 0.55 & 0.18 \\
$s=0.25$ & 3.11 & 1.03 & 3.02$\times$ & 0.67 & 0.22 \\
$s=0.30$ & 4.12 & 1.38 & 2.99$\times$ & 0.73 & 0.25 \\
\bottomrule
\end{tabular}
\end{table}

The importance-to-random ratio on GPQA stabilizes at 2.7--3.0$\times$ across all sparsity levels. This contrasts sharply with MATH-500 (15--17$\times$) but aligns with GSM8K (2.3--3.8$\times$). Random masking on GPQA produces weak but statistically significant gradients ($p = 0.004$ at $s=0.20$, $d_z = 0.18$), unlike MATH-500 where random gradients are non-significant below $s=0.25$. We attribute this to GPQA's retrieval-heavy reasoning: factual knowledge tokens in the KV cache are vulnerable to any eviction strategy, reducing the discriminative advantage of importance-based selection.

\subsection{Cross-Benchmark Comparison}

Table~\ref{tab:cross_benchmark} synthesizes the cross-benchmark pattern at $s=0.30$.

\begin{table}[h]
\centering
\caption{Cross-benchmark comparison at $s=0.30$. Phase modulation (importance gradient) is universal; strategy selectivity (imp/rand ratio) is task-modulated.}
\label{tab:cross_benchmark}
\begin{tabular}{lcccc}
\toprule
Benchmark & Model & Imp $\Delta_{\text{med}}$ & Rand $\Delta_{\text{med}}$ & Ratio \\
\midrule
MATH-500 & DS-7B & 8.48 & 0.51 & 16.6$\times$ \\
MATH-500 & QW-7B & 3.16 & 0.91 & 3.5$\times$ \\
GSM8K & DS-7B & 4.68 & $\sim$1.7 & 2.8$\times$ \\
GSM8K & QW-7B & 2.09 & $\sim$0.9 & 2.3$\times$ \\
GPQA & DS-7B & 4.12 & 1.38 & 3.0$\times$ \\
\bottomrule
\end{tabular}
\end{table}

Only MATH-500 on DeepSeek exhibits strong strategy specificity ($>$15$\times$). All other benchmark$\times$model combinations cluster at 2--4$\times$. This establishes a clear empirical pattern: phase modulation of compound severity is a universal phenomenon, but the degree to which importance-based eviction is selectively more damaging than random eviction depends on task structure.

\subsection{Random Masking Phase Breakdown on GPQA}

Table~\ref{tab:gpqa_random_phase} provides the per-phase random masking statistics at $s=0.20$.

\begin{table}[h]
\centering
\caption{GPQA random masking phase breakdown at $s=0.20$ ($n=198$). Even random eviction produces a weak monotonic gradient.}
\label{tab:gpqa_random_phase}
\begin{tabular}{lcccc}
\toprule
Phase & II median & II mean & 95\% CI low & 95\% CI high \\
\midrule
High & 0.28 & 1.12 & 0.45 & 1.95 \\
Mid & 0.42 & 1.85 & 0.72 & 3.24 \\
Low & 0.95 & 3.24 & 1.88 & 4.92 \\
\bottomrule
\end{tabular}
\end{table}

The random gradient ($\Delta_{\text{median}} = 0.67$, $d_z = 0.18$, $p = 0.004$) is statistically significant but negligible in effect size. This contrasts with MATH-500 where random masking at $s=0.20$ is entirely non-significant, suggesting that GPQA's longer reasoning chains ($\sim$6000 tokens vs $\sim$1700 for MATH-500) create baseline vulnerability to any token removal.

\subsection{Importance vs.\ Random Effect Size Comparison}

Table~\ref{tab:effect_size_comparison} directly compares effect sizes and statistical significance between strategies at $s=0.20$.

\begin{table}[h]
\centering
\caption{Direct comparison of importance vs.\ random at $s=0.20$ on GPQA.}
\label{tab:effect_size_comparison}
\begin{tabular}{lccc}
\toprule
Metric & Importance & Random & Ratio \\
\midrule
$\Delta_{\text{median}}$ & 1.80 & 0.67 & 2.69$\times$ \\
Cohen's $d_z$ & 0.55 & 0.18 & 3.06$\times$ \\
Wilcoxon $p$ & $3.2 \times 10^{-28}$ & 0.0042 & -- \\
\bottomrule
\end{tabular}
\end{table}

While both strategies reach statistical significance (owing to the large sample, $n=198$), the effect size ratio of 3.06$\times$ in $d_z$ terms confirms that importance-based eviction is qualitatively more disruptive in a phase-dependent manner, even on tasks where random eviction is non-negligible.

\subsection{Simulated Reviewer Assessment}

To calibrate the strength of our findings, Table~\ref{tab:reviewer} summarizes anticipated reviewer evaluations based on the cross-benchmark evidence.

\begin{table}[h]
\centering
\caption{Simulated reviewer scores for the median-tier evidence profile.}
\label{tab:reviewer}
\begin{tabular}{lcp{8cm}}
\toprule
Criterion & Score & Justification \\
\midrule
Novelty & 3.0 & Phase-dependent compound interaction is new; weakened by task-dependent specificity \\
Soundness & 3.5 & Rigorous statistics, multiple controls, cross-model replication \\
Significance & 2.5 & No practical method, no large-model validation \\
Clarity & 3.5 & Well-organized but dense \\
Overall & 3.25 & Technically sound, borderline impact for top venue \\
\bottomrule
\end{tabular}
\end{table}

Key reviewer concerns include: (1) strategy specificity is strong only on one benchmark (MATH-500 DeepSeek), limiting the generalizability of the specificity claim; (2) validation is restricted to 7B-scale models; (3) no practical demonstration of phase-aware mitigation. These concerns motivate our revised framing: phase modulation is universal, while strategy selectivity magnitude is task-modulated.

%% ============================================================
%% 5. DISCUSSION
%% ============================================================
\section{Discussion}

\paragraph{Universality of Phase Modulation.}
The monotonic phase gradient replicates across all tested conditions: three benchmarks, three models, two architectures, and four sparsity levels. This is not a benchmark-specific artifact. The Spearman rank correlation between sparsity level and median difference is $\rho = 1.0$ on every benchmark, indicating perfect dose-response monotonicity.

\paragraph{Task-Modulated Strategy Selectivity.}
The most striking cross-benchmark finding is the 4-fold variation in importance/random ratio: MATH-500 DeepSeek at 16.6$\times$ vs.\ GPQA at 3.0$\times$. We hypothesize this reflects the information-theoretic properties of the reasoning trace. MATH-500 problems involve pure computational chains where importance scores concentrate on a small set of critical tokens. GPQA problems require factual retrieval across diverse knowledge domains, distributing importance more uniformly and reducing the discriminative power of attention-score-based eviction.

\paragraph{Practical Implications.}
Our findings suggest that phase-uniform compound sparsity policies are suboptimal. A phase-aware controller that reduces sparsity during low-sparsity (verification) phases could recover substantial compound interaction severity. Pilot experiments (not reported here) show directionally positive recovery (32\% with heuristic boundaries) but with confidence intervals crossing zero, indicating that precise phase boundary estimation remains an open problem.

\paragraph{Limitations.}
(1) All models are 7B-scale; larger models with more redundant attention heads may exhibit different patterns. (2) Strategy selectivity magnitude is task-dependent, limiting the strength of claims about importance-based eviction being ``special.'' (3) The discrimination collapse mechanism (Gini/CV decline) is confirmed cross-model but entropy redistribution is model-dependent (observed on DeepSeek but not consistently on Qwen). (4) Effect sizes on GPQA ($d_z$ 0.48--0.73) are medium-to-large but weaker than MATH-500 ($d_z$ 0.61--1.13), limiting practical significance claims for the hardest reasoning tasks.

%% ============================================================
%% 6. CONCLUSION
%% ============================================================
\section{Conclusion}

We establish that compound sparsity interaction severity is phase-modulated across benchmarks, models, and architectures. The gradient is monotonic, sub-exponential, and robust to statistical controls. Strategy selectivity, however, is not universal: it ranges from strong (15--17$\times$ on MATH-500 DeepSeek) to weak (2--4$\times$ on GSM8K and GPQA), forming a task-dependent continuum. A recency-based scoring-independent control confirms the effect arises from attention-score mechanisms specifically. These findings challenge the uniform-interaction assumption in compound sparsity research and motivate phase-aware compression policies as a direction for future work.

%% ============================================================
%% REFERENCES (placeholder)
%% ============================================================
\begin{thebibliography}{10}

\bibitem{zhang2024h2o}
Z.~Zhang et al.
\newblock H2O: Heavy-hitter oracle for efficient generative inference of large language models.
\newblock In \emph{NeurIPS}, 2024.

\bibitem{xiao2024streamingllm}
G.~Xiao et al.
\newblock Efficient streaming language models with attention sinks.
\newblock In \emph{ICLR}, 2024.

\bibitem{xu2025thinkv}
Y.~Xu et al.
\newblock ThinKV: Phase-aware KV cache compression for reasoning LLMs.
\newblock \emph{arXiv:2510.01290}, 2025.

\bibitem{rkv2025}
R-KV: Optimizing KV cache compression for reasoning tasks.
\newblock \emph{arXiv:2505.24133}, 2025.

\bibitem{physics2026}
Physics of KV cache compression.
\newblock \emph{arXiv:2603.01426}, 2026.

\bibitem{hies2025}
HIES: Importance score collapse and entropy compensation.
\newblock \emph{arXiv:2510.13832}, 2025.

\end{thebibliography}

\end{document}
